\chapter{Control Flow Instructions}
\label{ch:control-flow}

\section{Overview}

Control flow in SIRC-1 is built from effective-address instructions. The real CPU opcodes compute an address into an
address-register pair; jumps and calls are produced by targeting the program-counter pair \reg{p}.

The SIRC-1 provides these control-flow and address-calculation operations:

\begin{itemize}
	\item \mnemonic{LDEA} -- Load Effective Address into an address-register pair
	\item \mnemonic{LDEL} -- Load Effective Address and Link; also saves the return address in \reg{l}
	\item \mnemonic{BRAN} -- PC-relative assembler meta-instruction for \texttt{LDEA p, ...}
	\item \mnemonic{BRSR} -- PC-relative assembler meta-instruction for \texttt{LDEL p, ...}
	\item \mnemonic{RETS} -- Return meta-instruction for \texttt{LDEA p, (\#0, l)}
	\item \mnemonic{LJMP} -- Absolute-jump assembler meta-instruction for \texttt{LDEA p, ...}
	\item \mnemonic{LJSR} -- Absolute-call assembler meta-instruction for \texttt{LDEL p, ...}
\end{itemize}

All encoded control-flow instructions can be conditional. If the condition is false, the instruction is skipped and has
no side effects, as described in Chapter~\ref{ch:reading-instructions}.

\section{Legal Forms}

\begin{table}[H]
	\centering
	\scriptsize
	\begin{tabular}{llll}
		\toprule
		\textbf{Instruction} & \textbf{Syntax}                      & \textbf{Encoding}  & \textbf{Notes}                                 \\
		\midrule
		LDEA                 & \texttt{LDEA dest, (\#offset, src)}  & 0x18               & Immediate-offset effective address             \\
		LDEA                 & \texttt{LDEA dest, (rS, src)}        & 0x19               & Register-offset effective address              \\
		LDEA                 & \texttt{LDEA dest, -(\#offset, src)} & 0x1A               & Pre-decrement source, immediate offset         \\
		LDEA                 & \texttt{LDEA dest, -(rS, src)}       & 0x1B               & Pre-decrement source, register offset          \\
		LDEL                 & \texttt{LDEL dest, (\#offset, src)}  & 0x1C               & Immediate-offset effective address and link    \\
		LDEL                 & \texttt{LDEL dest, (rS, src)}        & 0x1D               & Register-offset effective address and link     \\
		LDEL                 & \texttt{LDEL dest, (\#offset, src)+} & 0x1E               & Post-increment source, immediate offset        \\
		LDEL                 & \texttt{LDEL dest, (rS, src)+}       & 0x1F               & Post-increment source, register offset         \\
		BRAN                 & \texttt{BRAN \#disp}                 & \texttt{LDEA} meta & Assembles to \texttt{LDEA p, (\#disp, p)}      \\
		BRAN                 & \texttt{BRAN @label}                 & \texttt{LDEA} meta & PC-relative linker-computed displacement       \\
		BRSR                 & \texttt{BRSR \#disp}                 & \texttt{LDEL} meta & Assembles to \texttt{LDEL p, (\#disp, p)}      \\
		BRSR                 & \texttt{BRSR @label}                 & \texttt{LDEL} meta & PC-relative linker-computed call               \\
		RETS                 & \texttt{RETS}                        & \texttt{LDEA} meta & Assembles to \texttt{LDEA p, (\#0, l)}         \\
		LJMP                 & \texttt{LJMP src}                    & \texttt{LDEA} meta & Assembles to \texttt{LDEA p, (\#0, src)}       \\
		LJMP                 & \texttt{LJMP src, \#offset}          & \texttt{LDEA} meta & Assembles to \texttt{LDEA p, (\#offset, src)}  \\
		LJMP                 & \texttt{LJMP src, rS}                & \texttt{LDEA} meta & Assembles to \texttt{LDEA p, (rS, src)}        \\
		LJSR                 & \texttt{LJSR src}                    & \texttt{LDEL} meta & Assembles to \texttt{LDEL p, (\#0, src)}       \\
		LJSR                 & \texttt{LJSR src, \#offset}          & \texttt{LDEL} meta & Assembles to \texttt{LDEL p, (\#offset, src)}  \\
		LJSR                 & \texttt{LJSR src, rS}                & \texttt{LDEL} meta & Assembles to \texttt{LDEL p, (rS, src)}        \\
		LJSR                 & \texttt{LJSR (\#offset, src)+}       & \texttt{LDEL} meta & Assembles to \texttt{LDEL p, (\#offset, src)+} \\
		LJSR                 & \texttt{LJSR (rS, src)+}             & \texttt{LDEL} meta & Assembles to \texttt{LDEL p, (rS, src)+}       \\
		\bottomrule
	\end{tabular}
	\caption{Legal Control-Flow and Effective-Address Forms}
	\label{tab:control-flow-legal-forms}
\end{table}

\noindent\textbf{Address registers:} \texttt{src} and \texttt{dest} are address-register pairs: \reg{l}, \reg{a},
\reg{s}, or \reg{p}. The branch meta-instructions always use \reg{p} as both the source and destination address-register
pair.

\noindent\textbf{Aliased register writes:} Instructions that write the same address-register pair through more than one
write-back path have architecturally undefined behavior. Avoid forms such as \texttt{LDEA a, -(\#0, a)} and
\texttt{LDEL l, (\#0, l)+}. The assembler rejects known hazardous aliased forms.

\noindent\textbf{Shift suffixes:} Control-flow and effective-address displacements are not shifted. Register-offset
forms do not accept shift suffixes in the assembler.

\section{Common Semantics}

\begin{table}[H]
	\centering
	\begin{tabular}{lp{9cm}}
		\toprule
		\textbf{Topic}      & \textbf{Definition}                                                                                                                                                                                                                                                                                                          \\
		\midrule
		Condition false     & The instruction is skipped and has no side effects. The program counter advances to the next instruction normally.                                                                                                                                                                                                           \\
		Address calculation & The selected source address-register high word supplies the segment. The target low word is the selected source low word plus the immediate or register displacement.                                                                                                                                                        \\
		Pre-decrement LDEA  & The source low word is decremented by one word before the displacement is applied. The decremented source low word is written back.                                                                                                                                                                                          \\
		Post-increment LDEL & The target address is calculated from the original source value. The source low word is incremented by one word in write-back.                                                                                                                                                                                               \\
		Subroutine link     & \mnemonic{LDEL}, \mnemonic{BRSR}, and \mnemonic{LJSR} save the normal next-instruction address in \reg{l}.                                                                                                                                                                                                                   \\
		Status flags        & Preserved. Control-flow instructions do not support status update override syntax.                                                                                                                                                                                                                                           \\
		Timing              & 6 cycles, plus any instruction-fetch wait states. These instructions do not perform a data-memory bus access.                                                                                                                                                                                                                \\
		Exceptions          & Effective-address calculation can raise segment-overflow faults when \texttt{SR.A} is set; see the notes below.                                                                                                                                                                                                              \\
		Privilege           & In protected mode, address-register-pair write-back is allowed only when every written high word would remain unchanged. If any high word would change, the instruction raises a privilege-violation fault instead of completing write-back. Direct writes to high address registers also raise a privilege-violation fault. \\
		Fault side effects  & Program-counter, destination address-register, source auto-update, and link-register updates occur after condition evaluation. If a fault aborts execution before write-back, those effects do not complete.                                                                                                                 \\
		\bottomrule
	\end{tabular}
	\caption{Common Control-Flow Instruction Semantics}
	\label{tab:control-flow-common-semantics}
\end{table}

\noindent Control-flow and effective-address instructions do not perform data-memory access and do not raise data bus,
data bus-protection, or invalid-opcode faults in documented forms. Effective-address calculation can raise a
segment-overflow fault when \texttt{SR.A} (Trap on Address Overflow) is set and the 16-bit low-word address calculation
wraps. Instruction fetch can still raise the normal fetch faults before the instruction executes. If trace mode was
enabled when the instruction began, an instruction-trace fault is raised after the instruction commits.

\section{Branch Meta-Instructions}

\begin{instructionbox}{BRAN -- Branch Meta-Instruction}
	\textbf{Assembles to:} \mnemonic{LDEA} with \reg{p} as the source and destination

	\textbf{Syntax:}
	\begin{lstlisting}
BRAN #disp                    ; p = p + disp
BRAN @label                   ; p = p + label offset
BRAN|cond #disp               ; Conditional PC-relative branch
\end{lstlisting}

	\textbf{Operands:} \texttt{\#disp} is a 16-bit PC-relative word displacement. \texttt{@label} is resolved by the
	linker to a PC-relative displacement from the next instruction address.

	\textbf{Operation:}
	\begin{verbatim}
; Assembles to LDEA p, (#disp, p)
if condition:
    p = p + displacement
\end{verbatim}

	\textbf{Description:}

	Branches relative to the current program counter. \mnemonic{BRAN} is an assembler meta-instruction, not a distinct CPU
	opcode. Use \mnemonic{LDEA p, (\#offset, src)} or \mnemonic{LDEA p, (rS, src)} for non-PC-relative computed jumps.

	\textbf{Write-back:} Inherits \mnemonic{LDEA} write-back to \reg{p}.

	\textbf{Flags:} Preserved.

	\textbf{Exceptions:} Same as \mnemonic{LDEA}; may raise segment-overflow or protected-mode privilege faults.

	\textbf{Condition codes:} If the condition is false, no branch occurs and status flags are preserved.

	\textbf{Timing:} Same as \mnemonic{LDEA}: 6 cycles, plus instruction-fetch wait states.

	\textbf{Privilege:} Same as \mnemonic{LDEA}. Protected-mode branches must preserve the high word of \reg{p}.

	\textbf{Examples:}
	\begin{lstlisting}
; Unconditional branch forward
BRAN #20                      ; Jump ahead 20 words

; Backward branch (loop)
loop:
    ; ... loop body ...
    BRAN|!= @loop             ; Jump back if not zero
\end{lstlisting}

\end{instructionbox}

\begin{instructionbox}{BRSR -- Branch to Subroutine Meta-Instruction}
	\textbf{Assembles to:} \mnemonic{LDEL} with \reg{p} as the source and destination

	\textbf{Syntax:}
	\begin{lstlisting}
BRSR #disp                    ; l = next PC; p = p + disp
BRSR @label                   ; l = next PC; p = p + label offset
BRSR|cond #disp               ; Conditional PC-relative call
\end{lstlisting}

	\textbf{Operands:} \texttt{\#disp} is a 16-bit PC-relative word displacement. \texttt{@label} is resolved by the
	linker to a PC-relative displacement from the next instruction address.

	\textbf{Operation:}
	\begin{verbatim}
; Assembles to LDEL p, (#disp, p)
if condition:
    l = address_of_next_instruction
    p = p + displacement
\end{verbatim}

	\textbf{Description:}

	Calls a PC-relative subroutine by saving the return address in \reg{l} and branching to the computed target. Return
	with \mnemonic{RETS}, which assembles to \texttt{LDEA p, (\#0, l)}.

	\textbf{Write-back:} Inherits \mnemonic{LDEL} write-back to \reg{l} and \reg{p}.

	\textbf{Flags:} Preserved.

	\textbf{Exceptions:} Same as \mnemonic{LDEL}; may raise segment-overflow or protected-mode privilege faults.

	\textbf{Condition codes:} If the condition is false, no call occurs, \reg{l} is not written, and status flags are
	preserved.

	\textbf{Timing:} Same as \mnemonic{LDEL}: 6 cycles, plus instruction-fetch wait states.

	\textbf{Privilege:} Same as \mnemonic{LDEL}. Protected-mode calls must preserve the high words written to \reg{l} and
	\reg{p}.

	\textbf{Examples:}
	\begin{lstlisting}
; PC-relative call
BRSR @function

; Return from the subroutine
RETS
\end{lstlisting}

\end{instructionbox}

\section{Return Meta-Instruction}

\begin{instructionbox}{RETS -- Return from Subroutine Meta-Instruction}
	\textbf{Assembles to:} \texttt{LDEA p, (\#0, l)}

	\textbf{Syntax:}
	\begin{lstlisting}
RETS                          ; p = l
RETS|cond                     ; Conditional return
\end{lstlisting}

	\textbf{Operands:} None. The source address-register pair is implicitly \reg{l}; the destination is implicitly \reg{p}.

	\textbf{Operation:}
	\begin{verbatim}
; Assembles to LDEA p, (#0, l)
if condition:
    p = l
\end{verbatim}

	\textbf{Description:}

	Returns from a subroutine by copying the link register to the program counter. Used after \mnemonic{LJSR} or
	\mnemonic{BRSR}.

	\textbf{Write-back:} Inherits \mnemonic{LDEA} write-back to \reg{p}.

	\textbf{Flags:} Preserved.

	\textbf{Exceptions:} Same as \mnemonic{LDEA}; may raise protected-mode privilege faults if the return would change
	\reg{p}'s high word.

	\textbf{Condition codes:} If the condition is false, no return occurs and status flags are preserved.

	\textbf{Timing:} Same as \mnemonic{LDEA}: 6 cycles, plus instruction-fetch wait states.

	\textbf{Privilege:} Available in protected mode and supervisor mode, subject to protected-mode address-register
	write-back restrictions.

	\textbf{Examples:}
	\begin{lstlisting}
my_function:
    ; ... function body ...
    RETS                      ; Return to caller
\end{lstlisting}

\end{instructionbox}

\section{Load Effective Address}

\begin{instructionbox}{LDEA -- Load Effective Address}
	\textbf{Opcodes:} 0x18 (Immediate), 0x19 (Register), 0x1A (Pre-Dec Immediate), 0x1B (Pre-Dec Register)

	\textbf{Syntax:}
	\begin{lstlisting}
LDEA dest, (#offset, src)     ; dest = src + offset
LDEA dest, (rS, src)          ; dest = src + rS
LDEA dest, -(#offset, src)    ; src -= 1; dest = src + offset
LDEA dest, -(rS, src)         ; src -= 1; dest = src + rS
LDEA|cond dest, (#offset, src)
\end{lstlisting}

	\textbf{Operands:} \texttt{dest} and \texttt{src} are address-register pairs. Immediate forms use a 16-bit
	displacement. Register forms use \texttt{rS} as a displacement. Pre-decrement forms also update \texttt{src}.

	\textbf{Operation:}
	\begin{verbatim}
if condition:
    if pre_decrement:
        src.low = src.low - 1
    dest.high = src.high
    dest.low  = src.low + displacement
\end{verbatim}

	\textbf{Description:}

	Computes an effective address and loads it into the destination address-register pair. This is useful for pointer
	arithmetic and for control flow when the destination is \reg{p}. Pre-decrement forms also update the source address
	register pair.

	\textbf{Write-back:} If the condition is true, writes the computed address to \texttt{dest}. Pre-decrement forms also
	write the decremented source address back to \texttt{src}.

	\textbf{Flags:} Preserved.

	\textbf{Exceptions:} May raise a segment-overflow fault during effective-address calculation when \texttt{SR.A} is set.
	In protected mode, raises a privilege-violation fault if write-back would change a written address-register high word.

	\textbf{Condition codes:} If the condition is false, no destination or source address register is written and status
	flags are preserved.

	\textbf{Timing:} 6 cycles, plus instruction-fetch wait states. No data-memory bus access is performed.

	\textbf{Privilege:} Available in protected mode and supervisor mode, subject to protected-mode address-register
	write-back restrictions.

	\textbf{Examples:}
	\begin{lstlisting}
; Pointer arithmetic
LDEA a, (#16, a)              ; a = a + 16

; Computed jump
LDEA p, (r1, a)               ; p = a + r1

; Pre-decrement source pointer and copy nearby address
LDEA a, -(#0, s)              ; s -= 1; a = s
\end{lstlisting}

\end{instructionbox}

\begin{instructionbox}{LDEL -- Load Effective Address and Link}
	\textbf{Opcodes:} 0x1C (Immediate), 0x1D (Register), 0x1E (Post-Inc Immediate), 0x1F (Post-Inc Register)

	\textbf{Syntax:}
	\begin{lstlisting}
LDEL dest, (#offset, src)     ; l = next PC; dest = src + offset
LDEL dest, (rS, src)          ; l = next PC; dest = src + rS
LDEL dest, (#offset, src)+    ; l = next PC; dest = src + offset; src += 1
LDEL dest, (rS, src)+         ; l = next PC; dest = src + rS; src += 1
LDEL|cond dest, (#offset, src)
\end{lstlisting}

	\textbf{Operands:} \texttt{dest} and \texttt{src} are address-register pairs. Immediate forms use a 16-bit
	displacement. Register forms use \texttt{rS} as a displacement. Post-increment forms also update \texttt{src}.

	\textbf{Operation:}
	\begin{verbatim}
if condition:
    l = address_of_next_instruction
    dest.high = src.high
    dest.low  = src.low + displacement
    if post_increment:
        src.low = src.low + 1
\end{verbatim}

	\textbf{Description:}

	Computes an effective address, saves the return address to \reg{l}, and writes the computed address to the destination
	address-register pair. It is the real link-writing primitive behind \mnemonic{BRSR} and \mnemonic{LJSR}.

	\textbf{Write-back:} If the condition is true, writes the return address to \reg{l} and the computed address to
	\texttt{dest}. Post-increment forms also write the incremented source address back to \texttt{src}.

	\textbf{Flags:} Preserved.

	\textbf{Exceptions:} May raise a segment-overflow fault during effective-address calculation when \texttt{SR.A} is set.
	In protected mode, raises a privilege-violation fault if write-back would change a written address-register high word.

	\textbf{Condition codes:} If the condition is false, no destination, link, or source address register is written and
	status flags are preserved.

	\textbf{Timing:} 6 cycles, plus instruction-fetch wait states. No data-memory bus access is performed.

	\textbf{Privilege:} Available in protected mode and supervisor mode, subject to protected-mode address-register
	write-back restrictions.

	\textbf{Examples:}
	\begin{lstlisting}
; Absolute call through a
LDEL p, (#0, a)               ; l = next PC; p = a

; Explicit destination form
LDEL a, (r1, s)               ; l = next PC; a = s + r1

; Post-increment source pointer after call target calculation
LDEL p, (#0, a)+              ; l = next PC; p = old a; a += 1
\end{lstlisting}

\end{instructionbox}

\section{Long Jump Meta-Instructions}

\begin{instructionbox}{LJMP -- Long Jump Meta-Instruction}
	\textbf{Assembles to:} \mnemonic{LDEA} with \reg{p} as the destination

	\textbf{Syntax:}
	\begin{lstlisting}
LJMP src                      ; p = src
LJMP src, #offset             ; p = src + offset
LJMP src, rS                  ; p = src + rS
LJMP|cond src                 ; Conditional long jump
\end{lstlisting}

	\textbf{Operands:} \texttt{src} is an address-register pair. Optional \texttt{\#offset} or \texttt{rS} operands select
	immediate or register displacement forms.

	\textbf{Operation:}
	\begin{verbatim}
; Assembles to LDEA p, ...
if condition:
    p = src + displacement
\end{verbatim}

	\textbf{Description:}

	Performs an absolute jump through an address register. \mnemonic{LJMP} is an assembler meta-instruction, not a separate
	CPU opcode.

	\textbf{Write-back:} Inherits \mnemonic{LDEA} write-back to \reg{p}.

	\textbf{Flags:} Preserved.

	\textbf{Exceptions:} Same as \mnemonic{LDEA}; may raise segment-overflow or protected-mode privilege faults.

	\textbf{Condition codes:} If the condition is false, no jump occurs and status flags are preserved.

	\textbf{Timing:} Same as \mnemonic{LDEA}: 6 cycles, plus instruction-fetch wait states.

	\textbf{Privilege:} Same as \mnemonic{LDEA}. Protected-mode jumps must preserve the high word of \reg{p}.

\end{instructionbox}

\begin{instructionbox}{LJSR -- Long Jump to Subroutine Meta-Instruction}
	\textbf{Assembles to:} \mnemonic{LDEL} with \reg{p} as the destination

	\textbf{Syntax:}
	\begin{lstlisting}
LJSR src                      ; l = next PC; p = src
LJSR src, #offset             ; l = next PC; p = src + offset
LJSR src, rS                  ; l = next PC; p = src + rS
LJSR (#offset, src)+          ; l = next PC; p = src + offset; src += 1
LJSR (rS, src)+               ; l = next PC; p = src + rS; src += 1
LJSR|cond src                 ; Conditional long call
\end{lstlisting}

	\textbf{Operands:} \texttt{src} is an address-register pair. Optional \texttt{\#offset} or \texttt{rS} operands select
	immediate or register displacement forms. Post-increment forms also update \texttt{src}.

	\textbf{Operation:}
	\begin{verbatim}
; Assembles to LDEL p, ...
if condition:
    l = address_of_next_instruction
    p = src + displacement
    if post_increment:
        src.low = src.low + 1
\end{verbatim}

	\textbf{Description:}

	Calls a subroutine through an address register. The return address is saved in \reg{l}; return with \mnemonic{RETS}.
	Post-increment forms advance \texttt{src} after computing the target. Because \mnemonic{LJSR} implicitly writes \reg{l}
	and \reg{p}, post-increment source registers cannot be \reg{l} or \reg{p}.

	\textbf{Write-back:} Inherits \mnemonic{LDEL} write-back to \reg{l}, \reg{p}, and any post-increment source register.

	\textbf{Flags:} Preserved.

	\textbf{Exceptions:} Same as \mnemonic{LDEL}; may raise segment-overflow or protected-mode privilege faults.

	\textbf{Condition codes:} If the condition is false, no call occurs, \reg{l} is not written, any post-increment source
	register is preserved, and status flags are preserved.

	\textbf{Timing:} Same as \mnemonic{LDEL}: 6 cycles, plus instruction-fetch wait states.

	\textbf{Privilege:} Same as \mnemonic{LDEL}. Protected-mode calls must preserve the high words written to \reg{l} and
	\reg{p}.

	\textbf{Examples:}
	\begin{lstlisting}
; Call function at address in a
LOAD ah, #0x0040
LOAD al, #0x0000              ; a = 0x00400000
LJSR a                        ; Call 0x00400000

; Call with offset
LJSR a, #function_offset      ; p = a + function_offset

; Call with register displacement
LOAD r1, #8
LJSR a, r1                    ; p = a + r1

; Call through a table entry and advance to the next entry
LJSR (#0, a)+                 ; p = old a; a += 1
LJSR (r1, a)+                 ; p = a + r1; a += 1
\end{lstlisting}

\end{instructionbox}

\section{Control Flow Patterns}

\subsection{If-Then-Else}
\begin{lstlisting}
; if (r1 > r2) { ... } else { ... }
CMPR r1, r2
BRAN|<= @else_part
; Then part
; ...
BRAN @end
else_part:
; Else part
; ...
end:
\end{lstlisting}

\subsection{While Loop}
\begin{lstlisting}
; while (r1 > 0) { ... }
while_start:
    CMPI r1, #0
    BRAN|<= @while_end
    ; Loop body
    ; ...
    BRAN @while_start
while_end:
\end{lstlisting}

\subsection{Switch/Jump Table}
\begin{lstlisting}
; Switch on r1 (0-3)
LOAD r2, r1                   ; r2 = r1
ADDI r2, #0, LSL #1           ; r2 = (r2 << 1) = r2 * 2 (word size)
LDEA p, (r2, p)               ; Jump via table

jump_table:
    DW case0 - jump_table
    DW case1 - jump_table
    DW case2 - jump_table
    DW case3 - jump_table
\end{lstlisting}
